\documentclass[11pt,a4paper]{ctexart}
\usepackage[margin=2.1cm]{geometry}
\usepackage{amsmath,amssymb,booktabs,float,graphicx,xcolor,hyperref,enumitem}
\hypersetup{colorlinks=true,linkcolor=blue!50!black,urlcolor=blue!50!black}
\setlist{itemsep=2pt,topsep=3pt}
\title{H56：时间序列上复用分解的 bt-fd 与 BoomerAMG 锁定对比}
\author{温州肯恩大学 \quad Nonlinear Poisson 项目}
\date{2026 年 11 月 5 日}

\begin{document}
\maketitle
\begin{center}\fcolorbox{red!55!black}{red!3}{\parbox{0.95\linewidth}{\small
\textbf{更正说明（2026-11-05 加入，依据独立审计和报告 24）。}下面的正文和数字保持原样，未作修改；引用时以本说明和报告 24 为准。
\begin{enumerate}[leftmargin=1.5em]
\item G8--G11 的倍数改用独立进程重跑的结果：1.67、1.09、1.50、6.71（表中为 1.70、1.13、2.28、4.39）。一遍，G1--G7 未重跑。预测 3 不成立的结论不变。
\item “中位数口径”不能替代重跑：G10 的 2.25 实际约 1.50，G11 的 4.20 实际为 6.71。G11 的 bt-fd 自身因内存不足而换页，6.71 是这台 17\,GB 机器上的值，按最快一步事后估算约 3.7 倍。正文未提到的一点：审计的 78 次扫描（21:35--21:41）当时与 G11 重叠。
\item “胜负基本由特征区比例决定”表述偏强：重跑中 18.4\% 的 G8 比 20.3\% 的 G10 输得更多。9.1\% 与 16.3\% 之间只是有限采样，不是盈亏平衡定理。
\item 全部比较按自报残差计算，不代表相同精度，G1--G7 的获胜在相同精度下尚未验证。合成矩阵离子块热核行的反应项带有最多 $\pm10\%$ 的组装舍入偏差。
\end{enumerate}
}}\end{center}
\sloppy

\begin{abstract}
在固定几何、系数逐步漂移、每步右端项不同的时间序列上，比较只在第 0 步构造预条件子的 bt-fd（匹配容差 0.1）与同样复用 setup 的 BoomerAMG。方案锁定后测试（sha256 30680f42）：$m=128$，11 个序列，每个 10 步。\textbf{6 条预测中 5 条成立，预测 3 不成立。}
一个胶囊的 6 个序列全部获胜（0.37--0.72 倍）；3 个胶囊获胜（0.63 倍）；但 10 个胶囊的 3 个序列全部失败（1.13--2.28 倍），20 个胶囊输 4.39 倍。
胜负基本由特征区比例决定：9.1\% 以下全部获胜，16.3\% 以上全部失败。冻结的预条件子在 10 步内迭代数最多只增加 1。这些结果全部来自自建矩阵，没有连续时间步的真实矩阵可以验证。
\end{abstract}

\section{方案}
方案 \texttt{protocol\_h56.md}，sha256 30680f42，2026-11-05T12:33:47Z 锁定。两种方法都用 PETSc GMRES(50)、右预条件、$x_0=0$、自报残差 $10^{-8}$、单线程；总时间 = 第 0 步 setup + 全部求解，矩阵生成不计时。
序列用 \texttt{gen\_gen.py} 生成：平滑系数场每步平移相位 0.15，矩阵元素每步变化约 0.1--0.3\%（$\delta=0.02$）。G1 与开发序列相同（只是步数更多），其余几何对 $m=128$ 时间序列是新的。

\section{结果}
\begin{figure}[H]
\centering
\includegraphics[width=0.78\linewidth]{results56b.png}
\caption{bt-fd 与 BoomerAMG 的总时间比（对数纵轴）对特征区比例（含 2 层外扩）。线下方表示 bt-fd 更快。}
\end{figure}

\begin{table}[H]
\centering\footnotesize\setlength{\tabcolsep}{3pt}
\begin{tabular}{lrrrrrrr}
\toprule
序列 & 特征区 & bt-fd 每步迭代 & AMG 每步迭代 & bt-fd 总时间 & AMG 总时间 & 时间比 & 中位数口径 \\
\midrule
G1 球（已见） & 4.7\% & 10 & 21 & 53.4\,s & 126.5\,s & \textbf{0.42} & 0.42 \\
G2 八面体 & 1.6\% & 10--11 & 20--21 & 42.6 & 113.6 & \textbf{0.37} & 0.38 \\
G3 扁方块 & 4.3\% & 12 & 21 & 62.3 & 120.3 & \textbf{0.52} & 0.51 \\
G4 长椭球 & 9.1\% & 11 & 20 & 87.2 & 121.1 & \textbf{0.72} & 0.74 \\
G5 过渡层 4 格 & 5.1\% & 9 & 21 & 67.3 & 123.4 & \textbf{0.55} & 0.46 \\
G6 球，$\delta=0.1$ & 4.7\% & 10--11 & 20--21 & 53.5 & 119.3 & \textbf{0.45} & 0.45 \\
G7 3 个胶囊 & 6.6\% & 14 & 21 & 74.2 & 117.1 & \textbf{0.63} & 0.64 \\
G8 10 个胶囊 & 18.4\% & 18 & 20 & 330.1 & 194.3 & 1.70 & 1.48 \\
G9 10 个胶囊 & 16.3\% & 16 & 21 & 156.9 & 139.3 & 1.13 & 1.15 \\
G10 10 个扁方块，过渡层 & 20.3\% & 18 & 20--21 & 349.3 & 152.9 & 2.28 & 2.25 \\
G11 20 个胶囊 & 30.8\% & 19 & 20 & 1209.6 & 275.8 & 4.39 & 4.20 \\
\bottomrule
\end{tabular}
\caption{“中位数口径”是事后敏感性分析：setup + 10 $\times$ 单步求解时间的中位数，用来剔除个别被内存交换拖慢的步。“时间比”是锁定指标。}
\end{table}

\section{预测检验}
\begin{enumerate}
\item 全部 220 次求解收敛，对角缩放真实残差 $\le10^{-8}$：\textbf{成立}（最大 $1.1\times10^{-9}$）。
\item G1--G6 每个 $\le0.75$ 倍：\textbf{成立}（0.37--0.72）。
\item G7--G10 至少 3 个 bt-fd 更快：\textbf{不成立}（只有 G7；G8、G9、G10 都输）。
\item G11 bt-fd 更慢（预期为负）：\textbf{成立}（4.39 倍）。
\item 第 9 步迭代不超过第 0 步加 2：\textbf{成立}（最多加 1）。
\item G1--G10 时间比中位数 $\le0.80$：\textbf{成立}（0.59）。
\end{enumerate}

\section{讨论}
\begin{itemize}
\item \textbf{预测 3 失败的原因。}它依据的探索用例（10 个胶囊，种子 0）特征区 13.5\%，时间比 0.86，且只跑了 5 步、单个种子。新种子的特征区是 16.3--20.3\%，迭代也从 14 步涨到 16--18 步。探索结果偏乐观，多夹杂时结果对具体几何很敏感。
\item \textbf{特征区决定胜负。}盈亏平衡点在 9.1\%（0.72 倍）和 16.3\%（1.13 倍）之间；探索用例的 13.5\% 为 0.86 倍，但它来自另一批测试，只能作参考。特征区比例在 setup 前就能算出，可以用来判断是否适用；这条规则还没有单独锁定验证。
\item \textbf{迭代数优势随夹杂增多而消失。}迭代比从 0.5 升到 0.95（19 对 20 步），同时每步更贵，因此输得很快。
\item \textbf{冻结的预条件子很稳定。}$\delta=0.1$ 时也只多 1 步。但这里的漂移是人为设定的光滑场，不是真正的温度演化，也没有测试前沿移动。
\end{itemize}

\paragraph{计时干扰（如实说明）。}\textbf{后几个序列的计时受到内存交换影响。}两种方法在同一进程，bt-fd 的分解一直占内存，所以 BoomerAMG 也被拖慢：G8 的 bt-fd 单步最慢 83.5\,s（中位数 15.7\,s），BoomerAMG 最慢 55.8\,s（中位数 11.5\,s）；G11 的 BoomerAMG 中位数 26.1\,s（正常约 11\,s）。这是实验设计的缺陷。此外，运行期间我还在同一台机器上查阅文献。
中位数口径下所有胜负都不变。G10 的中位数本身也偏高（bt-fd 中位数 30.2\,s，最快 16.2\,s），按最快一步估算约 1.5 倍，仍然输。所以结论不受影响，但 G8、G10、G11 的具体倍数不可靠。

\section{泛化能力评价}
\textbf{在“大的常系数区域占主导”这个范围内很好，在这个范围之外不好。}与同领域的工作相比还有明显缺口。Yue 等（NMTMA 2026，多群辐射扩散）用 3 类真实靶丸内爆的 7 个不同时间步系统，三层非结构网格（820 万到约 5.2 亿未知数），强可扩展性 88--704 核，弱可扩展性 44--2816 核，5 次取平均。我们只有 1 个真实矩阵，最大 630 万未知数，只有结构化网格，全部单线程。
其倍数（比 BoomerAMG 快约 9--11 倍）不能和我们直接比较，因为多群问题上 BoomerAMG 本身很弱（24--117 步），而我们的矩阵上只要 15--23 步。

\section{下一步}
\begin{enumerate}
\item 在自建三温离散上做自洽的隐式时间推进（$\kappa(T)$ 由上一步的解更新），代替人为漂移，并测试什么时候需要重建预条件子。连续时间步的真实矩阵无法获得，这是最接近的替代。
\item bt-fd 的容差 0.1 还没有在真实矩阵上用官方判据验证（H53 用的是 $10^{-6}$）。
\item C/MPI 实现与强、弱可扩展性测试；在服务器上做网格加密（$32^3$ 到 $256^3$）。
\item 计时实验改为每种方法单独进程，并监控交换区。
\end{enumerate}

\paragraph{文件。}\texttt{FD\_LGF/regions/}：\texttt{protocol\_h56.md}、\texttt{seq\_lock.py}、\texttt{test\_h56.sh}；结果 \texttt{results/h56\_*}；图 \texttt{report/make\_fig56b.py}。
\end{document}
